\chapter{Wheezing}

\section*{Definition}
Wheezing is a high-pitched whistling sound when you breathe, usually when you breathe out. It means your airways have narrowed.

\section*{What Happens in Your Body}
Wheezing happens when air is forced through narrowed breathing passages. The type and location of the wheeze can help your doctor find where the blockage is.

\section*{Causes}
\begin{itemize}
  \item Asthma
  \item Chronic obstructive pulmonary disease (COPD)
  \item Bronchitis
  \item Heart failure (cardiac asthma)
  \item Foreign body aspiration
\end{itemize}

\section*{When to See a Doctor}
See your doctor if:
\begin{itemize}
  \item Symptoms are severe or getting worse over time
  \item Wheezing is persistent or accompanied by difficulty breathing
  \item You have other concerning symptoms such as fever, weight loss, or bleeding
\end{itemize}

\section*{Related Symptoms}
\begin{itemize}
  \item Chest tightness
  \item Shortness of breath
  \item Cough
  \item Use of accessory muscles
  \item Prolonged expiration
\end{itemize}
\textbf{Important:} If this symptom persists, your doctor will check for serious underlying causes.


\section*{Self-Care / Home Management}
\begin{itemize}
  \item Maintain adequate hydration with water, warm tea, or broth to thin secretions.
  \item Use a humidifier or steam inhalation to soothe irritated airways and loosen mucus.
  \item Rest and avoid exposure to smoke, dust, and other respiratory irritants.
  \item Over-the-counter remedies such as saline nasal spray, cough drops, or honey (for adults) may provide symptomatic relief.
  \item Monitor for warning signs including high fever, difficulty breathing, chest pain, or worsening symptoms.
\end{itemize}

\section*{How Your Doctor Will Evaluate You}
\begin{itemize}
  \item History focuses on onset, duration, character of symptoms (productive vs dry cough, nasal discharge color), fever pattern, and exposure history.
  \item Physical examination includes vital signs, lung auscultation, nasal inspection, and oropharyngeal examination.
  \item Targeted testing may include rapid antigen/ PCR for respiratory viruses, chest X-ray, pulse oximetry, and sputum culture.
  \item Allergy testing or pulmonary function tests are reserved for chronic or recurrent cases.
\end{itemize}

\section*{Treatment Options}
\begin{itemize}
  \item Symptomatic management includes antitussives, expectorants, decongestants, antihistamines, and saline irrigation as appropriate.
  \item Bacterial infections require appropriate antibiotic therapy based on suspected pathogen and local resistance patterns.
  \item Inhaled bronchodilators or corticosteroids are used for reactive airway disease or asthma exacerbations.
  \item Referral to pulmonology or ENT is indicated for chronic, recurrent, or treatment-refractory cases.
\end{itemize}

\section*{References}
\begin{thebibliography}{99}
\bibitem{wheezing:ref1} Global Initiative for Asthma (GINA). ``Global Strategy for Asthma Management and Prevention.'' 2023 Update.
\bibitem{wheezing:ref2} Global Initiative for Chronic Obstructive Lung Disease (GOLD). ``Global Strategy for the Diagnosis, Management, and Prevention of COPD.'' 2023 Report.
\bibitem{wheezing:ref3} Pasterkamp H, Kraman SS, Wodicka GR. ``Respiratory sounds: Advances beyond the stethoscope.'' \textit{American Journal of Respiratory and Critical Care Medicine}, 1997;156(3):974--987.
\bibitem{wheezing:ref4} Saglani S, McKenzie SA. ``Wheezing in children: A review.'' \textit{Archives of Disease in Childhood}, 2008;93(3):236--241.
\bibitem{wheezing:ref5} Ducharme FM, Bhogal SK. ``The role of reliever therapy in acute asthma.'' \textit{BMJ}, 2013;346:f1398.
\bibitem{wheezing:ref6} Custovic A, Jackson DJ, Simpson A. ``Wheezing in preschool children: When and how to treat?'' \textit{The Lancet Respiratory Medicine}, 2018;6(7):490--492.
\end{thebibliography}
